\documentclass[17pt]{beamer}
%\documentclass{beamer}
\usetheme{Madrid}
\useoutertheme{noslidenum}
\setbeamertemplate{navigation symbols}{}
\usepackage{beamerthemesplit}
\usepackage{ae,aecompl}
\usepackage[scaled=.95]{helvet}
\usepackage[english]{babel}
\usepackage[latin1]{inputenc}
\usepackage[T1]{fontenc}
\usepackage{mathpazo,courier,euler}
\usepackage{listings}
\lstset{language=sh,
    basicstyle=\ttfamily\bfseries,
  showstringspaces=false,
  keywordstyle=\color{black}\bfseries}

% theme split
\usepackage{verbatim}
\newenvironment{colorverbatim}[1][]%
{%
\color{blue}
\verbatim
}%
{%
\endverbatim
}%

\definecolor{NavyBlue}{RGB}{0, 76, 153}
\setbeamercolor{structure}{fg=NavyBlue}
\author[FOSSEE]{}
\institute[IIT Bombay]{}
\date[]{}

% logo
\logo{\includegraphics[height=1.30cm]{../images/3t-logo.pdf}}
\logo{\includegraphics[height=0.6cm]{../images/fossee-logo.png}

\hspace{200pt}
\includegraphics[scale=0.5]{../images/fossee-logo.png}\\
\hspace{275pt}
\includegraphics[scale=0.9]{../images/3t-logo.pdf}}

\begin{document}

\sffamily \bfseries
\title
[Getting Started With For Loop]
{Getting Started With For Loop}
\author
[FOSSEE, IIT - Bombay]
{\small Talk to a Teacher\\{\color{blue}\url{http://spoken-tutorial.org}}\\National Mission on Education
 through ICT\\{\color{blue}\url{http://sakshat.ac.in}} \\ [0.8cm]Script by: Hardik Ghaghada \\ Narration by: \\
 [0.5cm]
{\small 15 May 2013}}

%%%%%%%%%%%%%%%%%%%%%%%%%%%%%%%%%%%%%%%%%%%%%%%%%%%%%%%%%%%%%%%%%%%%%%%
\begin{frame}
  \titlepage
\end{frame}
%%%%%%%%%%%%%%%%%%%%%%%%%%%%%%%%%%%%%%%%%%%%%%%%%%%%%%%%%%%%%%%%%%%%%%%
\begin{frame}
\frametitle{Objectives}
  At the end of this tutorial, you will be able to, 
\begin{itemize}
\item Understand indentation in Python
\item Use the for loop
\item Use range() function
\item Write blocks in Python shell
\item Write blocks in IPython interpreter
\end{itemize}
\end{frame}
%%%%%%%%%%%%%%%%%%%%%%%%%%%%%%%%%%%%%%%%%%%%%%%%%%%%%%%%%%%%%%%%%%%%%%%
\begin{frame}[fragile]
\frametitle{Whitespace in Python}
\begin{itemize}
\item Whitespace is significant
\item Blocks are conventionally indented using 4 spaces
\begin{verbatim}
     Block A
     Block A
         Block B
         Block B
     Block A
\end{verbatim}
\end{itemize}
\end{frame}
%%%%%%%%%%%%%%%%%%%%%%%%%%%%%%%%%%%%%%%%%%%%%%%%%%%%%%%%%%%%%%%%%%%%%%%
\begin{frame}[fragile]
\frametitle{Exercise 1}
  Write a \verb~for~ loop which iterates through a list of numbers and find
  the square root of each number.
  The numbers are,
\begin{verbatim}
   1369, 7225, 3364, 7056,
   5625, 729, 7056, 
   576, 2916
\end{verbatim}
\end{frame}
%%%%%%%%%%%%%%%%%%%%%%%%%%%%%%%%%%%%%%%%%%%%%%%%%%%%%%%%%%%%%%%%%%%%%%%
\begin{frame}[fragile]
\frametitle{Exercise 2}
  Print the square root of numbers in the list.
  Numbers are,
\begin{verbatim}
7225, 3268, 3364, 2966, 7056,
5625, 729, 5547, 7056, 576, 2916
\end{verbatim}
\emph{do it in the ipython interpreter}
\end{frame}
%%%%%%%%%%%%%%%%%%%%%%%%%%%%%%%%%%%%%%%%%%%%%%%%%%%%%%%%%%%%%%%%%%%%%%%
\begin{frame}[fragile]
\frametitle{Exercise 3}
  Find out the cube of all the numbers from 1 to 10. \\
  \emph{do it in the python shell}
\end{frame}
%%%%%%%%%%%%%%%%%%%%%%%%%%%%%%%%%%%%%%%%%%%%%%%%%%%%%%%%%%%%%%%%%%%%%%%
\begin{frame}
\frametitle{\verb~range()~ function}
    \begin{itemize}
    \item built in function in Python
    \item generates a list of integers
        \begin{itemize}
        \item \emph{syntax:} range([start,] stop[, step])
        \item \emph{example:}
            \begin{itemize}
            \item range(1, 20) - \emph{generates integers from 1 to 20}
            \item range(20) - \emph{generates integers from 0 to 20}
            \end{itemize}
        \end{itemize}
    \end{itemize}
\end{frame}
%%%%%%%%%%%%%%%%%%%%%%%%%%%%%%%%%%%%%%%%%%%%%%%%%%%%%%%%%%%%%%%%%%%%%%%
\begin{frame}
\frametitle{Exercise 4}
  Print all the odd numbers from 1 to 50.
\end{frame}
%%%%%%%%%%%%%%%%%%%%%%%%%%%%%%%%%%%%%%%%%%%%%%%%%%%%%%%%%%%%%%%%%%%%%%%
\begin{frame}
\frametitle{Summary}
  In this tutorial,we learnt to,
\begin{itemize}
\item Create blocks in python using ``for''
\item Indent the blocks of code
\item Iterate over a list using ``for'' loop
\item Use the ``range()'' function
\end{itemize}
\end{frame}
%%%%%%%%%%%%%%%%%%%%%%%%%%%%%%%%%%%%%%%%%%%%%%%%%%%%%%%%%%%%%%%%%%%%%%%
\begin{frame}
\frametitle{Evaluation}
\begin{enumerate}
    \item Indentation is not mandatory in Python
    \begin{itemize}
        \item True
        \item False
    \end{itemize}
    \vspace{8pt}
    \item Write a ``for'' loop to print the product of all natural numbers from 1 to 20
    \vspace{8pt}
    \item What will be the output of-range(1,5)
\end{enumerate}
\end{frame}
%%%%%%%%%%%%%%%%%%%%%%%%%%%%%%%%%%%%%%%%%%%%%%%%%%%%%%%%%%%%%%%%%%%%%%%
\begin{frame}
\frametitle{Solutions}
\begin{enumerate}
    \item False
    \vspace{8pt}
    \item y = 1\\
    for x in range(1,21):\\
    \hspace{12pt}
              y*=x\\
        print y 
    \vspace{8pt}      
    \item {[1,2,3,4]}
\end{enumerate}
\end{frame}
%%%%%%%%%%%%%%%%%%%%%%%%%%%%%%%%%%%%%%%%%%%%%%%%%%%%%%%%%%%%%%%%%%%%%%%
\begin{frame}
\frametitle{FOSSEE}
{\color{NavyBlue}Free and Open Source Software for \\Education} \\
\begin{itemize}
\item To enable the use of free and open source software tools
\item For more details, please visit: {\color{blue}\url{http://fossee.in/}}
\end{itemize}
\end{frame}
%%%%%%%%%%%%%%%%%%%%%%%%%%%%%%%%%%%%%%%%%%%%%%%%%%%%%%%%%%%%%%%%%%%%%%%
\begin{frame}
\frametitle{About the Spoken Tutorial Project}
\begin{itemize}
\item Watch the video available at {\color{blue}\url{http://spoken-tutorial.org /What\_is\_a\_Spoken\_Tutorial}} 
\item It summarises the Spoken Tutorial project 
\item If you do not have good bandwidth, you can download and watch it
\end{itemize}
\end{frame}
%%%%%%%%%%%%%%%%%%%%%%%%%%%%%%%%%%%%%%%%%%%%%%%%%%%%%%%%%%%%%%%%%%%%%%%
\begin{frame}
\frametitle{Spoken Tutorial Workshops}The Spoken Tutorial Project Team 
\begin{itemize}
\item Conducts workshops using spoken tutorials 
\item Gives certificates to those who pass an online test 
\item For more details, please write to \\ \hspace {0.5cm}{\color{blue}contact@spoken-tutorial.org}
\end{itemize}
\end{frame}
%%%%%%%%%%%%%%%%%%%%%%%%%%%%%%%%%%%%%%%%%%%%%%%%%%%%%%%%%%%%%%%%%%%%%%%
\begin{frame}
\frametitle{Acknowledgements}
\begin{itemize}
\item Spoken Tutorial Project is a part of the Talk to a Teacher  project 
\item Supported by the National Mission on Education through  ICT, MHRD, Government of India 
\item More information: \\{\color{blue}\url{http://spoken-tutorial.org/NMEICT-Intro}}
\end{itemize}
\end{frame}
%%%%%%%%%%%%%%%%%%%%%%%%%%%%%%%%%%%%%%%%%%%%%%%%%%%%%%%%%%%%%%%%%%%%%%%
\begin{frame}

  \begin{block}{}
  \begin{center}
  \textcolor{blue}{\Large THANK YOU!} 
  \end{center}
  \end{block}
\begin{block}{}
  \begin{center}
    For more Information, visit our website\\
    {\color{blue}\url{http://fossee.in/}}
  \end{center}  
  \end{block}
\end{frame}
%%%%%%%%%%%%%%%%%%%%%%%%%%%%%%%%%%%%%%%%%%%%%%%%%%%%%%%%%%%%%%%%%%%%%%%
\end{document}
